\documentclass[10pt]{article}
\usepackage[margin=20mm]{geometry}
\usepackage{amsmath,amssymb,fontspec,unicode-math,xurl,hyperref}
\IfFontExistsTF{DejaVu Serif}{\setmainfont{DejaVu Serif}}{\setmainfont{Latin Modern Roman}}
\setmathfont{latinmodern-math.otf}
\setlength{\parindent}{0pt}\setlength{\parskip}{6pt}\sloppy
\title{Initial Scientific Review of
Rethinking Addressing in Language Models via Contextualized Equivariant Positional Encoding}
\author{Assigned review 0243 / I25-43612}
\date{Initial review; scientific adjudication pending}
\begin{document}
\maketitle
Abstract: Two demonstrated failures affect the central theoretical contribution: the constructive proof of computational expressivity uses an incorrect triangular inverse and an invalid exact decoder, and the shift-invariance proposition is false under its shared-rotation assumptions. Both are Major. A released-code discrepancy in the parameter-efficient training filter is Minor, without evidence that the historical experiments used identical code. These findings do not refute the existence of another expressivity construction or the reported benchmark measurements.\par
\section*{Identity, scope, and evidence}
The reviewed original is the 32-page PMLR version of Zhu et al. [paper], including Appendices A–E. PDF and printed page numbers coincide.\par
All 69 preserved source-file hashes match the staged manifest. Static inspection uses VITA-Group/TAPE commit `466d9e61e57dd19e2fe23fcec295ee323561f6fd`; its constructed pinned-source archive and member mappings remain preserved. The arithmetic submodule at commit `c86a748bf4544877c35cb8c8198f5847f511c33a` was unavailable in staged material. No author code, training, or benchmark was executed.\par
The core chains audited are Eqs. 4–8 through Appendix B.1.1 to Theorem B.1/Theorem 4.1, and Eqs. 4–5 to Proposition B.9/4.2. The independent check independently checked original mathematical pages, the counterexamples below, and relevant static code. Reviewer-written calculations are preserved in and .\par
\section*{Problem 1: The expressivity construction fails on legal inputs}
Assessment: Major. Theorem 4.1 (p. 7), formalized as Theorem B.1 (pp. 18–20), claims a constant-depth, constant-width construction solving the \(\Pi(5)\) group word problem exactly. It explicitly depends on Lemma B.7 (p. 22), which defines \(M_{ij}=1\) for \(j\leq i\), including the diagonal, and asserts \[\prod_{j=1}^{i}(I-\xi_j\xi_j^\top)
 =I-\sum_{j=1}^{i}\big[(I+M\odot\Xi\Xi^\top)^{-1}\Xi\big]_j\xi_j^\top.\] Here \(\Xi\) has rows \(\xi_j^\top\); the right side is interpreted as outer products.\par
\subsection*{A counterexample satisfying the construction}
Use the permitted sequence \([\mathrm{bos}],[1\leftrightarrow2]\). The proof’s embedding has \[\xi_1=0,\qquad \xi_2=v=(1,-1,0,0,0,0)^\top,\qquad
 \Xi\Xi^\top=\operatorname{diag}(0,2).\] The printed inverse gives \(\eta_2=v/3\). After the third layer, \[E_{\mathrm{printed}}=I-\tfrac13 vv^\top,\quad
 E_{\mathrm{required}}=I-vv^\top,\quad
 E_{\mathrm{printed}}v=\tfrac13v\ne -v=E_{\mathrm{required}}v.\] Thus it is not the transposition matrix. This is an exact identity failure independent of precision, optimization, or a difficult input distribution.\par
The second layer uses this inverse (pp. 18–19), and the third invokes Lemma B.7 to identify Eq. 12 with the cumulative permutation (p. 20). The fourth-layer query and key give the score \(a^\top E b\), where \[a=(1,2,3,4,5,54.5)^\top,\qquad b=(1,2,3,4,5,-1)^\top.\] For the actual transposition the score is \(54-54.5=-1/2\); for the printed construction it is \(55-1/3-54.5=1/6\). The sign test in the proof changes on this legal nonidentity word.\par
\subsection*{The exact decoder also fails}
The final step on p. 20 uses \(h(z)=\min(1,\max(0,z))\) and asserts that a positive input becomes label 1. Even with repaired permutation matrices, an identity prefix of length \(i=3\) (bos followed by the same swap twice) has score \(1/2\) and decoder input \[z=\frac{e^{1/2}}{e^{1/2}+2}-\frac13\approx0.11852943,\qquad h(z)=z\ne1.\] The bos-only prefix gives input zero and output 0 instead of the stipulated identity label 1. The specification requires precise Boolean outputs at every prefix (pp. 17–18); an unstated threshold is not the printed construction.\par
\subsection*{Counterevidence, repair, and impact}
The recurrence in Lemma B.7 sums only over \(j<i\). Replacing the mask inside the inverse by its strictly lower-triangular part repairs this WY identity; our check verifies the repair on the legal example. Removing the diagonal from the causal attention mask elsewhere would be a different change. A decoder repair and a full uniform exact complexity construction still need proof.\par
These defects invalidate the supplied derivation of a headline expressivity result and warrant Major even if individual steps are readily repairable. They do not establish that no architecture in the family can solve the problem or that practical arithmetic accuracy is wrong. The broad family permits operations beyond the experimental parameterization; not all computational-precision and architectural-realizability questions are settled here.\par
\section*{Problem 2: Shared orthogonal equivariance does not imply phase invariance}
Assessment: Major. Proposition 4.2 (p. 7) and Proposition B.9 (pp. 23–24) assert invariance to uniform shifts of positional indices for RoPE or reweighted random Fourier features, assuming Eqs. 4–5. Those equations apply one common \(O\in O(R)\) to the last coordinate of every block. A shift produces block-dependent rotations \(O_{\Delta,m}\). The proof displays that dependence and then substitutes a single \(O\) in its induction. General left mixing of the \(ML\) rows in Eq. 8 does not commute with different block rotations.\par
\subsection*{Explicit permitted two-layer counterexample}
Set \(N=C=M=2\), \(L=B=1\), \(R=2\), \(x_1=(1,0)\), \(x_2=(1,1)\), and causal mask \(\left(\begin{smallmatrix}1&0\\1&1\end{smallmatrix}\right)\). Use unit-weight Fourier features \[e_{i,1}=(1,0),\qquad e_{i,2}=(\cos(\pi i/2),\sin(\pi i/2)).\] Layer 1 has zero queries/keys and zero token value map with a token residual. Tokens remain unchanged while Eq. 7 averages the available positional rows. In Eq. 8 take \(W_1=I\), \(\psi(x)=I\), and \[W_2=S=\begin{pmatrix}1&1\\0&0\end{pmatrix}.\] The first positional block becomes \(F_1=(1,1)\), \(F_2=(1/2,1/2)\), and the second block is zero. Shifting all indices by \(\Delta=1\) before these operations instead gives \(F'_1=(0,0)\), \(F'_2=(1/2,-1/2)\).\par
In layer 2 take \(W_Q=W_K=\operatorname{diag}(1,0)\), \(\phi\) the scalar identity on the within-block dot product, and \(W_V=\left(\begin{smallmatrix}0&1\\0&0\end{smallmatrix}\right)\). Token 2 has logits \((1,1/2)\) before the shift and \((0,1/2)\) afterwards. The scalar values are 0 and 1, so its first output coordinate changes from \[\frac{1}{1+e^{1/2}}=0.3775406688
 \quad\hbox{to}\quad
 \frac{e^{1/2}}{1+e^{1/2}}=0.6224593312.\] A token residual adds the same quantity in both cases.\par
All parameters are finite. Zero queries/keys and the residual are admitted by the paper’s identity/averaging lemmas. The first positional operation is left multiplication and commutes with every shared right rotation. The second attention uses within-block inner products. Permuting tokens and mask together permutes both operations. Hence Eqs. 4–5 hold, but shift invariance fails. The theorem does not exclude zero frequency or require identical frequencies across blocks.\par
\subsection*{Counterevidence and bounded repair}
Appendix C (p. 24) uses a smaller experimental parameterization: learned matrices mix attention heads and share weights across \(M,L\). With matched frequency schedules across heads this avoids the cross-frequency mixing above. Table 8’s small shift-ID differences are meaningful evidence for that tested architecture, but cannot prove the unrestricted proposition. No failure of that table is claimed.\par
A sufficient corrected hypothesis requires invariance/equivariance under the actual blockwise shift subgroup, or the stronger product group \(\prod_m O(R)\). Restricting contextualization maps to commute with those rotations makes the induction valid. This repairs the symmetry implication under changed assumptions, without proving generalization or repairing Problem 1. The original false proposition supports a central claimed stability property, so Major is appropriate.\par
\section*{Problem 3: Released fine-tuning freezes a stated trainable map}
Assessment: Minor; released-version scope. Section 3.3 (p. 6) says the post-attention positional linear map is trained. In the pinned source, `models/llama/adape.py:727` defines singular `post\_attention\_linear`; initialization at lines 848–852 zeros its weights. But `train\_llama.py:230-237` freezes parameters unless names contain `pe` or plural `post\_attention\_linears`. The additional default whitelist is `embed,norm` (lines 65–67).\par
The actual name `model.layers.0.post\_attention\_linear.weight` matches none of these strings. In the documented base-model fine-tuning route it remains frozen. At `adape.py:775-780`, the zero projection is added to the positional residual, so this direct attention-derived positional update stays zero. The distinct positional MLP remains trainable.\par
An explicit whitelist override or a nonzero loaded checkpoint changes part of this scenario; the documented launch supplies no such override. Correcting the singular/plural whitelist addresses this static discrepancy, but no training repair was run. Historical experiment/checkpoint identity is unestablished, so this supports a local implementation finding without invalidating benchmark tables.\par
\section*{Coverage, positive checks, and unresolved items}
Main definitions, both central proof chains, experimental descriptions, and Appendix C–E interpretation were inspected. Common-rotation invariance of left positional maps and inner-product attention is consistent with Eqs. 4–5. The fusion speedup is \(5.63/2.56=2.1992\), consistent with \(2.2\times\). The YaRN relative gain is \(33.92/32.82-1=3.35\%\), consistent with \(3.4\%\). Table 3’s FLOP increase is about \(13.87\%\); calling it negligible is a qualitative judgment, not an independently reproduced performance result.\par
Missing arithmetic code, the unavailable separate attachment route, historical checkpoint identity, full precision/complexity accounting, and empirical reproduction remain unresolved. A limited primary metadata check found arXiv v2 dated 21 August 2025 [arxiv]; that is not a byte-level comparison or verified correction. No claim that no erratum exists is made.\par
The overall assessment is Major, based on Problems 1–2; Problem 3 is Minor. These are initial scientific judgments pending human review, with no acceptance decision.\par
9 J. Zhu, P. Wang, R. Cai, J. D. Lee, P. Li, and Z. Wang. Rethinking Addressing in Language Models via Contextualized Equivariant Positional Encoding. PMLR 267 (2025), 80155–80186. https://proceedings.mlr.press/v267/zhu25t.html. VITA-Group. TAPE source repository, pinned commit recorded above. https://github.com/VITA-Group/TAPE. J. Zhu et al. arXiv:2501.00712, version-history metadata, accessed during this review. https://arxiv.org/abs/2501.00712.\par
\end{document}
